% Local auxiliaries to OpenAI, A two-dimensional area law from a global
% spectral gap, Section 11, geometry:initial-stars, lines 333--370,
% especially 361--370, and prop:two-families, lines 299--323;
% revision adc7f1241b42e322a6451854ab7e4b4c146bf78a.
% These local assertions do not give the full isolated-star or
% recursive-repair theorem.

\section{Connected components near an initial mark}
\label{sec:area_law_initial_sector_components}

Fix an arbitrary origin $o\in\mathbb R^2$, a finite endpoint set
$Z\subset\mathbb Z^2$, a natural width $C\ge2$ and an initial layer
$k_0\ge50{,}000{,}000$. Choose arbitrary valid pitch residues
$a_h,b_h$ modulo $2^{p_h-\ell_h}$ at every layer $h$.
Let $\mathcal I$ be the resulting actual initial identifiers, let
$X_i^0$ be their closed birth regions, put
$U_i=\operatorname{int}(X_i^0)$, and write
$\chi(i)\in\{0,1\}$ for their initial colors.
Let $k\ge k_0$, let $z\in F_{k,\ell_k}$ be an actual fine-cell index,
and let $v$ be any of the nine marks of its defining cell. Set
$\ell=\ell_k-5$ and $r=2^\ell/2=t_k/64$.

Use the eight triangles obtained by subdividing every side of the
square centered at $v$ with half-side $r$ at its midpoint. Their
origin is $o_{\rm S}=v-r(1,1)$, their exponent is $\ell$, and their
cell index is $(0,0)$. Write $\mathcal J$ for the perimeter positions,
$\tau$ for their cyclic successor, $P_s^{\rm S}$ for the closed
triangles and $b_s^+$ for their final base endpoints.
Let $\sigma:\mathcal J\to\mathcal I$ be the canonical assignment of
Definition~\ref{def:area_law_initial_sector_assignment}, obtained from
the concentric working fan of half-side $2r$.
Color position $s$ by $\chi(\sigma(s))$. Retain an adjacency between
distinct consecutive positions precisely when these colors agree.
Write $\mathcal R$ for the runs of this graph and $[s]$ for the run
containing $s$. Their closed regions are
$Y_R=\bigcup_{s\in R}P_s^{\rm S}$.
Let $\lambda:\mathcal R\to\mathcal I$ be the unique assignment of
Theorem~\ref{thm:area_law_initial_sector_run_assignment}; thus
$\lambda([s])=\sigma(s)$.
Define
\begin{align}
    A      & =\{s\in\mathcal J:\sigma(s)\ne\sigma(\tau(s))\}, \\
    L      & =\bigcup_{s\in A}[v,b_s^+],                      \\
    \Omega & =B_\infty(v,r)\setminus L.
    \label{eq:area_law_initial_cut_square}
\end{align}
These are the local initial regions used in
\cite[Section~11]{OpenAI2026AreaLaw}.

\begin{theorem}[Connected components of the initial cut square]
    \label{thm:area_law_initial_sector_components}
    \lean{TNLean.PEPS.AreaLaw.Geometry.initialRegion_angular_connectedComponents_equiv}
    \leanok
    \uses{def:area_law_initial_sector_assignment,
        def:area_law_initial_region_colors,def:area_law_fan_runs,
        def:area_law_cell_fan_successor}
    Write $\mathcal C(\Omega)$ for the connected components of
    $\Omega$ in its relative topology, and $C_\Omega(x)$ for the
    component containing $x\in\Omega$.
    There is a unique bijection
    $\Phi:\mathcal C(\Omega)\longrightarrow\mathcal R$ such that,
    for every $x\in\Omega$ and $R\in\mathcal R$,
    \begin{align}
        \Phi(C_\Omega(x))=R
         & \quad\Longleftrightarrow\quad x\in Y_R.
        \label{eq:area_law_initial_component_membership}
    \end{align}
    The change set $A$ may be empty; the center is then retained.
    Separated runs with the same initial identifier correspond to
    distinct connected components.
\end{theorem}
\begin{proof}\leanok
    \uses{def:area_law_cell_fan,
        thm:area_law_cell_fan_successor_endpoints,
        thm:area_law_initial_sector_colors,
        thm:area_law_fan_cut_component_correspondence}
    The smaller fan has center $v$, since its origin is
    $v-r(1,1)$ and its half-side is $r$.
    The successor endpoint theorem makes $s$ and $\tau(s)$ adjacent.
    The actual adjacent-sector color theorem gives
    $\chi(\sigma(s))=\chi(\sigma(\tau(s)))$ if and only if
    $\sigma(s)=\sigma(\tau(s))$.
    Thus the color-change indices are exactly $A$, and their radial
    union is precisely $L$.
    Apply Theorem~\ref{thm:area_law_fan_cut_component_correspondence}
    to the smaller fan and its derived coloring. Its cut square is
    the set $\Omega$ in (\ref{eq:area_law_initial_cut_square}), its graph
    runs are $\mathcal R$, and its closed run regions are $Y_R$. The theorem therefore gives the unique
    bijection with (\ref{eq:area_law_initial_component_membership}).
\end{proof}

\begin{theorem}[A retained run belongs to its initial open region]
    \label{thm:area_law_initial_run_open_region}
    \lean{TNLean.PEPS.AreaLaw.Geometry.initialRegion_sector_run_mem_initialOpenRegion}
    \leanok
    \uses{thm:area_law_initial_sector_run_assignment,
        def:area_law_initial_open_regions,
        def:area_law_cell_fan_successor}
    For every run $R\in\mathcal R$ and every point
    $x\in\Omega\cap Y_R$, one has $x\in U_{\lambda(R)}$.
    This includes an empty change set and does not identify
    separated runs assigned to the same initial identifier.
\end{theorem}
\begin{proof}\leanok
    \uses{thm:area_law_initial_sector_run_assignment,
        thm:area_law_initial_active_successor,
        def:area_law_initial_open_regions}
    The canonical run decomposition gives
    \begin{align}
        X_{\lambda(R)}^0\cap\overline B_\infty(v,r)
         & =\bigcup_{\substack{T\in\mathcal R \\\lambda(T)=\lambda(R)}}Y_T.
    \end{align}
    The term $T=R$ contains $x$, so $x\in X_{\lambda(R)}^0$.
    Suppose that $x$ also belongs to its frontier.
    Since $x\in B_\infty(v,r)$,
    Theorem~\ref{thm:area_law_initial_active_successor} gives an
    index $s\in A$ with $x\in[v,b_s^+]$. This places $x$ in $L$,
    contrary to $x\in\Omega$.
    Hence $x\in X_{\lambda(R)}^0\setminus\partial X_{\lambda(R)}^0
        =\operatorname{int}(X_{\lambda(R)}^0)=U_{\lambda(R)}$.
\end{proof}
